\documentclass[11pt]{article}
\usepackage[T1]{fontenc}
\usepackage[a4paper,margin=29mm]{geometry}
\usepackage{amsmath,amssymb,amsthm,mathtools,booktabs}
\usepackage{microtype}
\usepackage{needspace}
\usepackage[hidelinks]{hyperref}
\numberwithin{equation}{section}
\newtheorem{theorem}{Theorem}[section]
\newtheorem{proposition}[theorem]{Proposition}
\newtheorem{lemma}[theorem]{Lemma}
\newtheorem{corollary}[theorem]{Corollary}
\theoremstyle{definition}
\newtheorem{definition}[theorem]{Definition}
\theoremstyle{remark}
\newtheorem{remark}[theorem]{Remark}
\newcommand{\C}{\mathbb C}
\newcommand{\Q}{\mathbb Q}
\newcommand{\R}{\mathbb R}
\newcommand{\Z}{\mathbb Z}
\newcommand{\Fp}{\mathbb F_p}
\newcommand{\dd}{\,\mathrm d}
\newcommand{\h}{\mathsf h}
\newcommand{\CT}{\operatorname{CT}}
\newcommand{\SL}{\operatorname{SL}}
\newcommand{\GL}{\operatorname{GL}}
\newcommand{\SU}{\operatorname{SU}}
\newcommand{\Sym}{\operatorname{Sym}}
\newcommand{\Supp}{\operatorname{Supp}}
\newcommand{\conv}{\operatorname{conv}}
\newcommand{\Lie}{\operatorname{Lie}}
\newcommand{\ord}{\operatorname{ord}}
\newcommand{\inv}{\operatorname{inv}}
\DeclareMathOperator{\Cart}{\mathcal C}
\setlength{\emergencystretch}{2em}
\raggedbottom
\title{Counterexamples to the Mathieu Property\\[6pt]
{\large From Explicit Quotient Constructions\\
to a General Classification}}
\author{Roy van Rijn}
\date{17 November 2026}
\hypersetup{pdftitle={Counterexamples to the Mathieu Property: From Explicit Quotient Constructions to a General Classification},
            pdfauthor={Roy van Rijn}}
\begin{document}
\maketitle

\begin{abstract}
We construct regular functions $P,Q$ on each of $\SL_2/2O$ and
$\SL_2/2I$ for which every positive Haar moment of $P$ vanishes, whereas
$\h(QP^m)$ is nonzero for infinitely many $m$. The same elementary
construction also supplies a radial witness on the defining module of
$\SU(2)$. These three examples give the following classification: for a
complex homogeneous space $G/H$, with $G$ and $H$ having reductive identity
components, the kernel of normalized Reynolds integration is
Mathieu--Zhao exactly when $[G^0,G^0]\subseteq H$. We give the reduction
to the examples, including disconnected acting groups. The counterexamples
are proved directly from Haar integration, polynomial identities and
Dirichlet's theorem. For either exceptional example and each positive
degree, one explicit fixed multiplier, chosen independently of the
complex coefficients, gives a transcendental ordinary Cauchy transform
for every polynomial recomposition of that degree. Over an arbitrary coefficient algebra,
we also determine the pure and safe elements of the polynomial subalgebra
generated by the example.
\end{abstract}

\section{The problem and the results}\label{sec:intro}

Let $A$ be a commutative complex algebra and let $\h:A\to\C$ be a
linear functional with $\h(1)=1$. A function $P\in A$ is called
\emph{pure} if
\[
   \h(P^m)=0\qquad(m\geq1).
\]
We call $P$ \emph{safe} if, for each fixed $Q\in A$, the moments
$\h(QP^m)$ vanish for all sufficiently large $m$. The cutoff is allowed
to depend on $P$ and $Q$. Thus $\ker\h$ is a Mathieu--Zhao subspace
precisely when every pure function is safe. This definition is equivalent
to the version using eventual vanishing of the pure moments: take a
sufficiently large power of $P$ and then treat the finitely many residue
classes of exponents separately.

For example, let $A=\C[t,t^{-1}]$ and let $\h=\CT$ extract the
constant term. The function $P=t$ is pure and safe: for
$Q=\sum_jq_jt^j$, one has $\CT(Qt^m)=q_{-m}$, which is zero once
$m$ exceeds the finite range of negative exponents of $Q$. The question
is whether this eventual vanishing persists beyond tori.

Mathieu's Main Conjecture asks whether every pure $K$-finite function
on a compact connected Lie group $K$ is safe
\cite[pp.~264--265]{Mathieu}. Here $K$-finite means that the translates
span a finite-dimensional vector space. Long gave an explicit
counterexample on $\SU(2)$ \cite[Theorem~4.2]{LongSU2}, then classified
the compact connected group case: the Mathieu property holds exactly
for tori \cite[Theorem~1.1]{LongCompact}. Regular functions on a complex
reductive group restrict to $K$-finite functions on a compact form.
Thus the case $H=\{1\}$ of the classification below is already covered
by Long's compact-group theorem.

We are interested in $A=\C[G/H]$, where $G$ is a complex linear algebraic
group, $H\subseteq G$ is a closed algebraic subgroup, and both identity
components $G^0$ and $H^0$ are reductive. The algebraic Matsushima
criterion gives affineness of $G^0/H^0$
\cite[Proposition~2.2]{Arzhantsev}. Each component of $G/H$ is a finite
quotient of a conjugate of this affine variety, so $G/H$ is affine too.
Its normalized Reynolds functional is the unique $G$-invariant functional
with value one on the constant function one. When $G$ is connected, it
is integration over a compact form, with total mass one.

For connected $G$, the space $G/H$ is \emph{spherical} if a Borel
subgroup has an open orbit. For reductive $H$ in this setting, Long
proves that the Mathieu property and equality of the pure set with
the nullcone are both equivalent to $[G,G]\subseteq H$
\cite[Theorem~1.1]{LongSpherical}. We extend the Mathieu-property classification
to arbitrary reductive homogeneous spaces and disconnected acting groups.

\begin{theorem}\label{thm:classification}
Let $G$ and $H$ be as above. Then
\[
 \ker\h\subseteq\C[G/H]\text{ is Mathieu--Zhao}
 \quad\Longleftrightarrow\quad [G^0,G^0]\subseteq H.
\]
If the containment fails, there are fixed regular functions $P,Q$ on
$G/H$ such that $P$ is pure and $\h(QP^m)\ne0$ for infinitely many $m$.
\end{theorem}

The important word in the second assertion is \emph{fixed}. A different
multiplier at each order would not disprove the Mathieu property. Our
examples use one polynomial $Q$, independent of $m$.

A counterexample on $\SU(2)$ does not automatically descend to
$\SL_2/\Gamma$, with $\Gamma\subseteq\SU(2)$ finite: both functions must be invariant under right
multiplication by $\Gamma$. Averaging $P$ enforces invariance, but
powers of that average contain mixed products of translates of $P$,
whose moments are not controlled by $\h(P^m)=0$. Our binary-form
construction retains all-power vanishing while supplying the required
right invariance. The exceptional quotients below have dimension three,
while a Borel subgroup of $\SL_2$ has dimension two; they supply the
nonspherical cases left by the rank-one reduction.

Here are the two exceptional examples explicitly. Write
\[
 M=\begin{pmatrix}x&y\\z&w\end{pmatrix},\qquad xw-yz=1,
\]
and take one of the following binary forms:
\begin{equation}\label{eq:binaryforms}
 \begin{aligned}
 f_8(X,Y)&=X^8+14X^4Y^4+Y^8,\\
 f_{12}(X,Y)&=XY(X^{10}+11X^5Y^5-Y^{10}).
 \end{aligned}
\end{equation}
The groups $2O$ and $2I$ are the binary octahedral and binary icosahedral
subgroups of $\SU(2)$. Their standard embeddings fix $f_8$ and $f_{12}$,
respectively. An elementary verification is given in
Appendix~\ref{app:invariance}.

For $n=8$ or $12$, put
\begin{equation}\label{eq:construction}
 \begin{gathered}
 b=\frac n2,\qquad a=b-1,\qquad U=f_n(x,y),\qquad F=f_n(z,w),\\
 R=\frac{f_n(z+Fx,w+Fy)}{F},\qquad P=U^aR^b,\qquad Q=F^a.
 \end{gathered}
\end{equation}
Although $R$ is written as a quotient, it is a polynomial with integer
coefficients. We will prove this before using it.

The parameters used throughout the construction are collected below.
The first row also gives the radial example in Section~\ref{sec:radial}.
Here $c=b(n-1)$ and $p=cm-a$ is the prime index for the exceptional
examples; $N(T)$ denotes the normalizer of the diagonal torus.
\begin{center}
\begin{tabular}{@{}lrrrrl@{}}
\toprule
Binary form & $n$ & $a$ & $b$ & $c$ & Quotient \\
\midrule
$X^2Y^2$ & 4 & 1 & 2 & 6 & $\SL_2/N(T)$ \\
$f_8$ & 8 & 3 & 4 & 28 & $\SL_2/2O$ \\
$f_{12}$ & 12 & 5 & 6 & 66 & $\SL_2/2I$ \\
\bottomrule
\end{tabular}
\end{center}

\begin{theorem}\label{thm:exceptional}
The functions in \eqref{eq:construction} are regular invariants on
$\SL_2/2O$ when $n=8$, and on $\SL_2/2I$ when $n=12$. In both cases
\[
 \h(P^m)=0\qquad(m\geq1).
\]
For the same fixed $Q$, one has
\begin{equation}\label{eq:primevaluation}
 \begin{array}{ll}
 v_p\bigl(\h(QP^m)\bigr)=-1&\text{if }n=8\text{ and }p=28m-3
 \text{ is prime},\\
 v_p\bigl(\h(QP^m)\bigr)=-1&\text{if }n=12\text{ and }p=66m-5
 \text{ is prime}.
 \end{array}
\end{equation}
Consequently both exceptional quotients fail the Mathieu property.
\end{theorem}

Here $v_p$ is the usual valuation on rational numbers, normalized by
$v_p(p)=1$. The moments in this theorem are rational. Dirichlet's theorem \cite[Theorem~18.1]{Sutherland}
gives infinitely many primes in the progressions $25\pmod{28}$ and
$61\pmod{66}$, so \eqref{eq:primevaluation} proves nonvanishing at
arbitrarily large orders.

The proof has two independent parts. Torus weights and integration by
parts give purity. For the marked moments, a congruence reduces Haar
integration to one coefficient of a binary form modulo $p$. Two explicit
Klein identities give a map from the curve defined by the binary form
to a nonsingular cubic curve. A differential identity and compatibility
of the Cartier operator with pullback transfer the coefficient
calculation to that cubic. There it reduces to a binomial coefficient
that is nonzero modulo $p$.

For the classification, equivariant pullback carries a bad pair to
any homogeneous space mapping onto one of these examples. A simple
factor of rank at least two instead produces a two-dimensional
$\SL_2$-summand in a representation, to which the radial example applies.
In the remaining case, inactive higher-rank factors can be removed,
leaving only rank-one factors. A proper stabilizer projection gives a
map onto one of the three quotients in the table; otherwise a diagonal
coupling again produces such a summand. If $[G,G]\subseteq H$ for
connected $G$, the quotient is a torus. For disconnected
$G$, a finite-trace argument separates the moment conditions on its
torus components.

Sections~\ref{sec:haar}--\ref{sec:klein} prove
Theorem~\ref{thm:exceptional}. Section~\ref{sec:radial} gives the radial
example, and Sections~\ref{sec:connected}--\ref{sec:components} prove
Theorem~\ref{thm:classification}. Section~\ref{sec:components} also
distinguishes purity from the invariant-theoretic nullcone. The final
section treats polynomial recompositions, ordinary Cauchy transforms
and coefficient algebras with nilpotents. We use standard structural
facts about reductive groups, the classical finite subgroups of $\SL_2$,
the two branching rules stated in Lemma~\ref{lem:doublet}, and the
Duistermaat--van der Kallen theorem in its explicit form
\eqref{eq:dvk}. All special moment identities used in the proof are
derived here.

\section{Haar integration and a polynomial construction}\label{sec:haar}

On $\SU(2)$ the second row is determined by the first:
\begin{equation}\label{eq:su2coords}
 z=-\overline y,\qquad w=\overline x,\qquad |x|^2+|y|^2=1.
\end{equation}
Normalized Haar measure is normalized surface measure on this unit
three-sphere. The variables $|x|^2$, $\arg x$ and $\arg y$ have product
measure
\[
 \dd t\,\frac{\dd\theta}{2\pi}\,\frac{\dd\phi}{2\pi},
 \qquad 0\leq t\leq1.
\]
One can see this directly by writing
$x=\sqrt t\,e^{i\theta}$ and $y=\sqrt{1-t}\,e^{i\phi}$ in the surface
element. In particular, for nonnegative integers $i,k$,
\begin{equation}\label{eq:beta}
 \h(x^iy^kz^kw^i)
 =(-1)^k\int_0^1t^i(1-t)^k\dd t
 =(-1)^k\frac{i!k!}{(i+k+1)!}.
\end{equation}
Every monomial whose two phase exponents do not cancel has integral zero.

We shall also use left invariance. The compact diagonal subgroup gives
weight $+1$ to $x,y$ and weight $-1$ to $z,w$; every nonzero weight has
zero integral. The complexified infinitesimal left action contains
\[
 D=z\partial_x+w\partial_y.
\]
Thus $\h(DT)=0$ for every regular polynomial $T$, and the product rule gives
\begin{equation}\label{eq:ibp}
 \h(SD^jT)=(-1)^j\h((D^jS)T).
\end{equation}
These identities involve polynomials on $\SL_2$ only.

We first work with any homogeneous $f\in\Z[X,Y]$ of even degree $n=2b\geq4$.
Set $a=b-1$ and define $U,F,R,P,Q$ as in
\eqref{eq:construction}, with $f$ in place of $f_n$.
Taylor expansion in a new scalar $t$ gives
\[
 f(z+tx,w+ty)=\sum_{k=0}^n H_k t^k,\qquad H_0=F.
\]
Substitution $t=F$ shows
\begin{equation}\label{eq:Rregular}
 R=1+\sum_{k=1}^n F^{k-1}H_k\in\Z[x,y,z,w].
\end{equation}
If a subgroup fixes $f$, its simultaneous right action on the two rows
fixes all the polynomials in this formula. Hence the construction respects
right invariants.

\begin{proposition}\label{prop:pure}
For every such binary form $f$, the polynomial $P$ is pure.
\end{proposition}
\begin{proof}
Fix $m\geq1$, and write $A=am$ and $B=bm$. Define
\[
 H_k^{(B)}=[t^k]f(z+tx,w+ty)^B\qquad(0\leq k\leq nB).
\]
Expanding $R^B$ gives
\begin{equation}\label{eq:weightdecomposition}
 P^m=U^A\sum_{k=0}^{nB}F^{k-B}H_k^{(B)}.
\end{equation}
Every summand is a polynomial, including those with $k<B$.
Indeed, at most $k$ of the $B$ factors contributing to $H_k^{(B)}$ can
use a positive power of $t$. The remaining factors each contribute $F$,
so $F^{B-k}$ divides $H_k^{(B)}$ over $\Z$.

The $k$th summand has left weight
\[
 nA-(n-2)k=(n-2)(B-k).
\]
Only $k=B$ can contribute to its Haar integral. Homogeneity also gives
\begin{equation}\label{eq:reversecoefficient}
 H_k^{(B)}=\frac{D^{nB-k}(U^B)}{(nB-k)!}.
\end{equation}
To verify this identity, replace $t$ by $t^{-1}$ in
$f(x+tz,y+tw)^B$ and multiply by $t^{nB}$.
Consequently, with $j=(n-1)B$,
\[
 \h(P^m)=\frac1{j!}\h\bigl(U^AD^j(U^B)\bigr).
\]
The degree of $U^A$ in the top row is $nA=(n-2)B<j$.
Equation~\eqref{eq:ibp} therefore makes the last expression zero.
\end{proof}

\begin{proposition}\label{prop:markedcomponent}
Put $c=b(n-1)$. For $m\geq1$, set $A=am$, $B=bm$ and $s=b(m-1)$.
Then
\begin{equation}\label{eq:T_m}
 \h(QP^m)=\h(T_m),\qquad
 T_m=\frac{U^A}{F}[t^s]f(z+tx,w+ty)^B.
\end{equation}
The displayed $T_m$ is an integer polynomial, homogeneous of degree
\begin{equation}\label{eq:rowdegree}
 N=cm-b
\end{equation}
in each row. In particular its Haar integral has denominator dividing
$(cm-a)!$.
\end{proposition}
\begin{proof}
Multiplication by $Q=F^a$ changes the weight in
\eqref{eq:weightdecomposition} to $(n-2)(B-k-b)$.
The zero weight is therefore $k=s=B-b$, and its power of $F$ is $-1$.
The same factor-counting argument shows that $F^b$ divides the coefficient
in \eqref{eq:T_m}, so division by $F$ is harmless.
The top degree is $nA+s$ and the bottom degree is $nB-s-n$;
both equal $cm-b$. Formula~\eqref{eq:beta} gives the denominator bound.
\end{proof}

\section{A congruence for the marked moments}\label{sec:congruence}

\begin{lemma}\label{lem:haarcongruence}
If $T\in\Z[x,y,z,w]$ is homogeneous of degree $p-1$ in each row,
where $p$ is an odd prime, then $p\h(T)$ is $p$-integral and
\begin{equation}\label{eq:haarcongruence}
 p\h(T)\equiv[X^{p-1}Y^{p-1}]T(X,Y,X,Y)\pmod p.
\end{equation}
\end{lemma}
\begin{proof}
It suffices to consider $x^iy^{p-1-i}z^kw^{p-1-k}$. Its integral
vanishes unless $i+k=p-1$. In that case \eqref{eq:beta} gives
$(-1)^ki!k!/p!$. Wilson's theorem implies
\[
 i!(p-1-i)!\equiv(-1)^{i+1}\pmod p,
\]
so multiplication by $p$ makes this integral congruent to one.
The same condition $i+k=p-1$ selects precisely the monomials on the
right of \eqref{eq:haarcongruence}.
\end{proof}

Choose an index $m$ for which
\begin{equation}\label{eq:primeindex}
 p=cm-a\quad\text{is prime}.
\end{equation}
Then $N=p-1$ in Proposition~\ref{prop:markedcomponent}. At two equal rows,
\[
 [t^s]f(X+tX,Y+tY)^B=\binom{nB}{s}f(X,Y)^B.
\]
It follows that
\begin{equation}\label{eq:keycongruence}
 p\h(QP^m)\equiv\binom{nB}{s}\mathcal H_p\pmod p,
 \qquad
 \mathcal H_p=[X^{p-1}Y^{p-1}]f(X,Y)^{(p-1)/b}.
\end{equation}
We used $A+B-1=(p-1)/b$ here. Notice also that
\[
 nB=p+q,\qquad q=bm+a=s+n-1,\qquad 0\leq s\leq q<p.
\]
Taking the coefficient of $t^s$ in
$(1+t)^{p+q}\equiv(1+t^p)(1+t)^q$ proves
\begin{equation}\label{eq:binomialunit}
 \binom{nB}{s}\equiv\binom qs\not\equiv0\pmod p.
\end{equation}
Thus the entire nonvanishing problem has become the calculation of
one binary coefficient $\mathcal H_p$.

\section{The two binary coefficients}\label{sec:klein}

For each form in \eqref{eq:binaryforms}, define the companion $G$ by
\begin{equation}\label{eq:companions}
 \begin{aligned}
 G_8(X,Y)&=X^2Y^2(X^4-Y^4)^2,\\
 G_{12}(X,Y)&=X^{20}+Y^{20}-228X^{15}Y^5
                  +494X^{10}Y^{10}+228X^5Y^{15}.
 \end{aligned}
\end{equation}
Let $v=\deg G$, so $v=12$ or $20$, and put
\[
 J=\frac{f_XG_Y-f_YG_X}{nv}.
\]
The polynomial identities we need are
\begin{equation}\label{eq:kleinrelations}
 \begin{array}{ll}
 J^2=f^3G/36-3G^3,&n=8,\\
 J^2=12f^5+G^3/144,&n=12.
 \end{array}
\end{equation}
These identities can be checked by substitution and multiplication;
the definitions specify every coefficient. They are a convenient
normalization of the classical Klein identities. Appendix~\ref{app:identities}
gives explicit verifications in scalar coordinates.

We use a coefficient-extraction rule for differentials in characteristic
$p$. Its short algebraic proof is included in
Appendix~\ref{app:cartier}. If $L$ is a one-variable function field over
$\Fp$ and $t$ is a separating variable, every differential has a unique
expression
\[
 \omega=\sum_{i=0}^{p-1}a_i^p t^i\dd t.
\]
Here ``separating'' means that $L$ is finite and separable over $\Fp(t)$.
The rule $\Cart(\omega)=a_{p-1}\dd t$ is independent of the chosen
separating variable and commutes with separable field extensions.
It is usually called the Cartier operator. In particular,
\begin{equation}\label{eq:cartierrule}
 \Cart\left(\frac{\sum_jc_jt^j}{Y^p}\dd t\right)
 =\frac1Y\sum_{r\geq0}c_{pr+p-1}t^r\dd t
 \quad(c_j\in\Fp).
\end{equation}

Work now over $\Fp$, with $p>5$ and $p\equiv1\pmod b$, and form
the field $L=\Fp(t,Y)$ defined by
\[
 Y^b=f(t,1).
\]
The polynomials $f(t,1)$ have simple roots for these primes: for $f_8$
the relevant quadratic discriminant is $14^2-4=192$, and for $f_{12}$
it is $11^2+4=125$. Together with $p\nmid4\cdot5$, these facts check
squarefreeness directly. The polynomial $Y^b-f(t,1)$ is Eisenstein at
any simple root of $f$ after extending constants; hence it defines a
field and $t$ is separating.

In $L$ put
\[
 g=\frac{G(t,1)}{Y^{v/2}},\qquad
 j=\frac{J(t,1)}{Y^{(n+v-2)/2}}.
\]
Equation~\eqref{eq:kleinrelations} gives
\begin{equation}\label{eq:elliptic}
 \begin{array}{ll}
 j^2=g/36-3g^3,&n=8,\\
 j^2=12+g^3/144,&n=12.
 \end{array}
\end{equation}
These are nonsingular cubic equations for $p>5$. More importantly,
differentiation gives the exact identity
\begin{equation}\label{eq:differentialidentity}
 \frac{\dd g}{j}=-v\frac{\dd t}{Y}.
\end{equation}
Indeed, Euler's homogeneous formula yields
$J=f'G/n-fG'/v$ on the chart with second coordinate one, while
$Y'/Y=f'/(bf)$. Substituting these into the derivative of $g$ proves
\eqref{eq:differentialidentity}. Since $p>5$ and $v\in\{12,20\}$,
the scalar $-v$ lies in $\Fp^\times$. Thus $\dd g\ne0$.
Over the perfect field $\Fp$, this means that $g$ is separating:
$L/\Fp(g)$ is separable. The cubic function field
$E=\Fp(g,j)$ defined by \eqref{eq:elliptic} is an intermediate field,
so the inclusion $E\subseteq L$ is separable as well.

Homogeneity identifies the binary coefficient with the affine one:
\[
 \mathcal H_p=[t^{p-1}]f(t,1)^{(p-1)/b},
\]
because $f(X,Y)^{(p-1)/b}$ has total degree $2(p-1)$.
Since $Y^{p-1}=f(t,1)^{(p-1)/b}$, the extraction rule gives
\[
 \Cart\left(\frac{\dd t}{Y}\right)
 =\mathcal H_p\frac{\dd t}{Y}.
\]
The numerator has degree at most $2(p-1)$, so only the coefficient
of $t^{p-1}$ survives. If the cubic on the right of
\eqref{eq:elliptic} is denoted by $h(g)$, the same calculation gives
\[
 \Cart\left(\frac{\dd g}{j}\right)
 =[g^{p-1}]h(g)^{(p-1)/2}\frac{\dd g}{j}.
\]
Applying $\Cart$ to \eqref{eq:differentialidentity} and using separable
pullback compares these two eigenvalues. Cartier is $\Fp$-linear, so
the common nonzero factor $-v\in\Fp^\times$ cancels. The two
coefficients are therefore equal. They are
\begin{equation}\label{eq:hassevalues}
 \begin{aligned}
 n=8:\quad
 \mathcal H_p&=\binom{(p-1)/2}{(p-1)/4}
                     (-1/12)^{(p-1)/4},\\
 n=12:\quad
 \mathcal H_p&=\binom{(p-1)/2}{(p-1)/3}
                     12^{(p-1)/6}144^{-(p-1)/3}.
 \end{aligned}
\end{equation}
Every factor is nonzero in $\Fp$. For primes in
\eqref{eq:primeindex}, the required congruence $p\equiv1\pmod b$
is automatic. Equations~\eqref{eq:keycongruence} and
\eqref{eq:binomialunit} now show that $p\h(QP^m)$ is a $p$-adic unit.
This proves \eqref{eq:primevaluation}. Proposition~\ref{prop:pure},
Appendix~\ref{app:invariance} and Dirichlet's theorem complete the proof
of Theorem~\ref{thm:exceptional}.

\section{A radial witness from the same construction}\label{sec:radial}

Long's compact-group argument already supplies a radial counterexample
for invariant measures on $\C^2$ \cite[Theorem~4.1]{LongCompact}.
We give an alternative pair and its closed marked-moment formula.
Take $f_4(X,Y)=X^2Y^2$ in Section~\ref{sec:haar}. The resulting pair
has the particularly simple form
\begin{equation}\label{eq:normalizerpair}
 P_4=x^2y^2(1+xzw^2)^4(1+yz^2w)^4,\qquad Q_4=z^2w^2.
\end{equation}
Indeed, the intermediate polynomial $R$ is
$(1+xzw^2)^2(1+yz^2w)^2$. Both functions in
\eqref{eq:normalizerpair} are invariant under the right diagonal torus
and under
\[
 \begin{pmatrix}0&-1\\1&0\end{pmatrix}.
\]
The latter transformation interchanges the two parenthesized factors.
Thus they descend to $\SL_2/N(T)$, where $N(T)$ is the normalizer of
the diagonal torus.

The first moment illustrates the cancellation. On $\SU(2)$, put
$t=|x|^2$ and $u=\overline{xy}$. Then
\[
 P_4=(xy)^2(1-tu)^4(1+(1-t)u)^4,\qquad Q_4=u^2.
\]
Phase integration keeps the coefficient of $u^2$ for $P_4$ and the
constant coefficient for $Q_4P_4$. Consequently,
\[
 \h(P_4)=\int_0^1t^2(1-t)^2(6-28t+28t^2)\dd t=0,
 \qquad \h(Q_4P_4)=\int_0^1t^2(1-t)^2\dd t=\frac1{30}.
\]
The following calculation extends this distinction to every order.

\begin{lemma}\label{lem:normalizerpositive}
For every $m\geq1$,
\begin{equation}\label{eq:normalizermoments}
 \h(P_4^m)=0,\qquad
 \h(Q_4P_4^m)=\beta_m
 =\frac{2m}{3m+1}\frac{(2m)!(4m)!}{(6m)!}>0.
\end{equation}
\end{lemma}
\begin{proof}
The phase calculation above also proves purity at every order, directly
from the binomial theorem:
\[
 \begin{aligned}
 \h(P_4^m)
 &=\frac1{(6m+1)!}\sum_{r=0}^{2m}(-1)^r
   \binom{4m}{r}\binom{4m}{2m-r}(2m+r)!(4m-r)!\\
 &=\frac{((4m)!)^2}{(6m+1)!(2m)!}
       \sum_{r=0}^{2m}(-1)^r\binom{2m}{r}=0.
 \end{aligned}
\]
For the marked moment, Proposition~\ref{prop:markedcomponent} gives
\[
 T_m=(xy)^{2m}(zw)^{-2}
       [t^{2m-2}](z+tx)^{4m}(w+ty)^{4m}.
\]
Set $L=2m-2$. Expanding this coefficient and using
\eqref{eq:beta} gives
\begin{equation}\label{eq:radialsum}
 \beta_m=\frac{((4m)!)^2}{(6m-1)!L!}
 \sum_{r=0}^L\frac{(-1)^r\binom Lr}
 {(4m-r)(4m-r-1)(2m+r+1)(2m+r+2)}.
\end{equation}
There is no division by $zw$ in the polynomial $T_m$; its displayed
coefficient is divisible by $(zw)^4$.

Let $S_m$ denote the sum in \eqref{eq:radialsum}. Use
$1/(k(k+1))=\int_0^1u^{k-1}(1-u)\dd u$ for each pair of consecutive
denominator factors and sum the binomial expansion. This gives
\[
 S_m=\int_0^1\!\int_0^1
     s^{2m}t^{2m}(1-s)(1-t)(t-s)^L\dd s\dd t>0.
\]
The positivity is immediate because $L=2m-2$ is even. For its value,
use symmetry to integrate over $0<s<t<1$ and put $s=rt$. With
\[
 I_m=\int_0^1r^{2m}(1-r)^L\dd r
     =\frac{(2m)!L!}{(4m-1)!},
\]
the $t$-integral and then the two $r$-integrals give
\[
 S_m=2I_m\left(\frac1{6m(6m+1)}
       -\frac{2m+1}{4m(6m+1)(6m+2)}\right)
     =\frac{I_m}{6m(6m+2)}.
\]
Substitution into \eqref{eq:radialsum} gives
\eqref{eq:normalizermoments}.
\end{proof}

The same pair works for much more than surface measure. For
$v=(v_0,v_1)\in\C^2$, define
\[
 M(v)=\begin{pmatrix}v_0&-\overline{v_1}\\v_1&\overline{v_0}\end{pmatrix},
 \qquad \mathcal P(v)=P_4(M(v)),\quad
 \mathcal Q(v)=Q_4(M(v)).
\]
These are polynomials in the real and imaginary coordinates of $v$.
They are unchanged by $v\mapsto e^{i\theta}v$, because
\eqref{eq:normalizerpair} is right-torus invariant.

\begin{proposition}[Radial identity]\label{prop:radial}
For every compactly supported finite positive $\SU(2)$-invariant measure
$\mu$ on $\C^2$ and every $m\geq1$,
\begin{equation}\label{eq:radialidentity}
 \int\mathcal P^m\dd\mu=0,\qquad
 \int\mathcal Q\mathcal P^m\dd\mu
 =\beta_m\int\|v\|^{12m-4}\dd\mu.
\end{equation}
The second integral is positive if $\mu$ is not concentrated at zero.
\end{proposition}
\begin{proof}
The matrix identity $M(kv)=kM(v)$ holds for $k\in\SU(2)$.
For $v\ne0$, put $\rho=\|v\|$ and write $M(v)=\rho g$ with $g\in\SU(2)$.
Using $t=|x|^2$ and $u=\overline{xy}$ for the unit-radius coordinates
of $g$, we have
\[
 P_4(\rho g)=\rho^4(xy)^2(1-\rho^4tu)^4(1+\rho^4(1-t)u)^4,
 \qquad Q_4(\rho g)=\rho^4u^2.
\]
Average the integrands first over $k$.
Phase integration selects order $2m$ for the pure moment and $2m-2$
for the marked moment. The averages are therefore $\rho^{12m}\h(P_4^m)=0$
and $\rho^{12m-4}\beta_m$, respectively. The formulas hold at $v=0$ as well.
Integrating these averaged identities against $\mu$ proves the result.
\end{proof}

\section{Connected reductive groups}\label{sec:connected}

We now prove Theorem~\ref{thm:classification} for connected $G$.
We use compact forms as in \cite[Section~3]{AdamsTaibi} and the
central-isogeny and complete-reducibility facts for reductive groups
\cite[Propositions~14.2 and~14.10, Sections~24.1 and~24.4]{Borel}.
The root-doublet and subdirect-product arguments follow Long
\cite[Proposition~3.1 and Lemmas~4.1, 4.2, 6.2]{LongSpherical}.
We give the reductions here and use
Proposition~\ref{prop:radial} for the moment input.

We use two elementary facts about normalized Reynolds functionals.
If $q:X\to Y$ is a surjective equivariant morphism of homogeneous spaces,
then
\begin{equation}\label{eq:pullback}
 \h_X(q^*f)=\h_Y(f).
\end{equation}
Both sides are normalized invariant functionals on $\C[Y]$, and such a
functional is unique. The same observation applies when the acting group
maps onto the acting group of $Y$. Also, a central isogeny
$\widetilde G\to G$ identifies
$\widetilde G/\pi^{-1}H$ with $G/H$ and preserves normalized integration.
Thus all witnesses constructed after this change are functions on the
original homogeneous space.

We call a submodule a \emph{defining doublet} if it is isomorphic to the
standard two-dimensional $\SL_2$-module.

\begin{proposition}[A line and a doublet]\label{prop:linedoublet}
Suppose an irreducible $G$-module $V$ has a nonzero $H$-stable line
$\C\xi$, and some homomorphism $\SL_2\to G$ has a defining
two-dimensional submodule $W\subseteq V$. Then $\C[G/H]$ contains
fixed functions $P,Q$ with
\[
 \h(P^m)=0,\qquad \h(QP^m)>0\qquad(m\geq1).
\]
\end{proposition}
\begin{proof}
Choose a compact form $L$ of $H$ and a compact form $K$ of $G$ containing
$L$. This is possible also for disconnected $H$; see
\cite[Theorems~3.7--3.8]{AdamsTaibi}. By conjugating the $\SL_2$
homomorphism and $W$ together, arrange that $\SU(2)$ maps into $K$.
Give $V$ a $K$-invariant Hermitian form. Orthogonal projection $\Pi_W$
commutes with $\SU(2)$, since both $W$ and $W^\perp$ are stable.
The map
\[
 \Phi:K\to W,\qquad \Phi(k)=\Pi_W(k\xi)
\]
is $\SU(2)$-equivariant. The pushforward of Haar measure is therefore
a compact positive invariant measure on $W$.
It is not concentrated at zero: continuity and full support of Haar
measure would otherwise give $K\xi\subseteq W^\perp$,
whereas the span of $K\xi$ is all of the irreducible module $V$.

Identify $W$ with $\C^2$ by an $\SU(2)$-equivariant unitary isomorphism
and apply Proposition~\ref{prop:radial} to this measure. The polynomials
$P=\mathcal P\circ\Phi$ and $Q=\mathcal Q\circ\Phi$ are regular on $G$:
their holomorphic factors are matrix coefficients. In an orthonormal
basis, $\overline{\rho(k)}=(\rho(k)^{-1})^{\mathsf T}$ for $k\in K$,
so the conjugate factors extend as matrix coefficients of the dual
representation. Moreover $\Phi(kl)=\chi(l)\Phi(k)$ for the character
$\chi$ of the $H$-stable line. On $L$ that character is unitary.
The phase invariance of $\mathcal P,\mathcal Q$ therefore makes $P,Q$
right $L$-invariant. Zariski density of $L$ in $H$ makes their regular
extensions right $H$-invariant. Thus they belong to $\C[G/H]$ and
\eqref{eq:radialidentity} proves the claim.
\end{proof}

\begin{lemma}\label{lem:doublet}
Every nontrivial finite-dimensional rational representation of a connected
simple complex group of rank at least two contains a defining doublet
for some root $\SL_2$.
\end{lemma}
\begin{proof}
For type $A_2$, an irreducible module has highest weight $(a,b)$,
$a,b\geq0$. The $\GL_3\downarrow\GL_2$ interlacing rule says that
the polynomial highest weight $(a+b,b,0)$ contains every pair satisfying
\[
 a+b\geq\mu_1\geq b\geq\mu_2\geq0.
\]
If $b>0$, use $(\mu_1,\mu_2)=(b,b-1)$; otherwise use $(1,0)$.
In the latter case nontriviality gives $a\geq1$. The difference is one,
so the root $\SL_2$ has a defining summand.

For $B_2$, write a highest weight as $(a,b)$ with $a\geq b\geq0$,
both integral or both half-integral. The
$\operatorname{Spin}_5\downarrow\operatorname{Spin}_4$ rule contains
all $(c,d)$ of the same parity with
\[
 a\geq c\geq b\geq|d|.
\]
The two $\SL_2$ highest weights are $c-d,c+d$.
Use $(c,d)=(b,b-1)$ when $b>0$, and $(1,0)$ otherwise.
The first choice is valid even when $b=\tfrac12$, since
$|b-1|\leq b$; the second uses $a\geq1$.
These are the classical branching rules; see
\cite[Theorem~2.1 and equation~(4.4)]{Molev}.

Every simple root system of rank at least two contains a root-generated
subalgebra of type $A_2$ or $B_2$: take adjacent simple roots, using
the long-root $A_2$ inside $G_2$ if necessary. A nontrivial representation
of a simple Lie algebra is faithful, so its restriction to this subalgebra
is nontrivial. A nontrivial irreducible constituent and the preceding
branching arguments give the doublet. Root homomorphisms integrate it
to the group, also after composing a simply connected cover.
\end{proof}

Call a normal simple factor \emph{active} if it acts nontrivially on
$G/H$. Since $G/H$ is affine, an active factor acts nontrivially on some
irreducible submodule
$U\subseteq\C[G/H]$. Evaluation at $eH$ is a nonzero $H$-fixed vector
of $U^*$. It is nonzero because a $G$-stable space of functions vanishing
at $eH$ would vanish at every point. Thus Lemma~\ref{lem:doublet} and
Proposition~\ref{prop:linedoublet} handle every active factor of rank
at least two.

\begin{lemma}\label{lem:subdirect}
Let $G=Z\times\SL_2^r$, with $Z$ a torus, and let $H$ be reductive.
If every projection $H\to\SL_2$ is surjective, then either
$[G,G]\subseteq H$ or the hypothesis of
Proposition~\ref{prop:linedoublet} holds.
\end{lemma}
\begin{proof}
The projections of $H^0$ are surjective as well, since their images have
finite index in the connected target. The center of $\Lie(H^0)$ maps
to zero in each $\mathfrak{sl}_2$. Hence its derived algebra projects
onto every factor. Decompose that semisimple algebra into simple ideals.
For each coordinate, exactly one ideal maps nontrivially: its image is
an ideal of $\mathfrak{sl}_2$, and two nonzero such images could not
commute. That ideal maps isomorphically onto $\mathfrak{sl}_2$.
Every simple ideal appears in at least one coordinate, since it cannot
map nontrivially to the torus. Thus the coordinates partition into
blocks on which the derived algebra is diagonal, up to inner
automorphisms. After conjugating individual factors,
\[
 [H^0,H^0]=\prod_B\Delta_B\SL_2.
\]
The equality holds for groups because connected algebraic subgroups
in characteristic zero are determined by their Lie algebras. The
connected center of $H^0$ lies in $Z$.
If every block is a singleton, $[G,G]\subseteq H$. Otherwise choose
two indices $i,j$ in one block. On the irreducible module
$V=\C_i^2\otimes\C_j^2$, that block acts diagonally and the connected
center acts trivially. Hence the entire fixed subspace is
\[
 V^{H^0}=\C\bigl(e_{i,1}\otimes e_{j,2}-e_{i,2}\otimes e_{j,1}\bigr),
\]
where the bases correspond under the chosen diagonal identifications.
This is the alternating summand in
$\C^2\otimes\C^2=\Sym^2\C^2\oplus\bigwedge^2\C^2$.
Normality of $H^0$ in $H$ makes this whole fixed subspace, and thus
its unique line, $H$-stable. The $i$th factor has the defining
submodule $\C_i^2\otimes\C e_1$ in $V$.
\end{proof}

\begin{proof}[Proof of Theorem~\ref{thm:classification} for connected $G$]
If $[G,G]\subseteq H$, then $G/H$ is a torus and its normalized
functional is constant term: $H$ is the inverse image of a closed
subgroup of the torus $G/[G,G]$, hence is normal with torus quotient.
We use the Duistermaat--van der Kallen
theorem \cite[pp.~222--224]{DvK} in the following form:
\begin{equation}\label{eq:dvk}
 \CT(f^m)=0\ (m\geq1),\ f\ne0
 \quad\Longrightarrow\quad0\notin\conv(\Supp f).
\end{equation}
A strictly positive linear functional on the finite support of $f$
eventually dominates the weights of any fixed multiplier $q$.
Therefore $\CT(qf^m)=0$ for all large $m$.

Now assume $[G,G]\not\subseteq H$. Pass to a finite central cover
$Z\times\prod S_i\to G$. If an active factor has rank at least two,
the preceding argument supplies the required pair. Otherwise quotient
out the higher-rank factors contained in the stabilizer. They are normal
and act trivially, so this leaves an isomorphic homogeneous space with
acting group $Z\times\SL_2^r$.

Consider the image $H_i$ of the stabilizer in one $\SL_2$ factor.
This image is closed and reductive. If it is proper and positive-dimensional,
its identity component is a torus; the full image normalizes that torus
and so lies in a conjugate of $N(T)$. If finite, the classical finite-subgroup classification
\cite[p.~215]{Flicker} puts it in one of the following groups:
\[
 N(T),\qquad 2O,\qquad 2I.
\]
Here cyclic and binary dihedral groups lie in $N(T)$, while $2T$ lies
in $2O$. In each case projection gives a genuine surjection of the
original homogeneous space onto the corresponding quotient of $\SL_2$.
Pull back the pair from Lemma~\ref{lem:normalizerpositive} or
Theorem~\ref{thm:exceptional}, using \eqref{eq:pullback}.
Finally, if every projection is surjective, apply
Lemma~\ref{lem:subdirect} and Proposition~\ref{prop:linedoublet}.
The alternative $[G,G]\subseteq H$ is excluded. This proves the
connected classification and its fixed-multiplier assertion.
\end{proof}

\section{Disconnected groups and moment nullcones}\label{sec:components}

We first give the finite-trace fact needed for disconnected spaces.
It does not rely on the signs of complex moments.

\begin{lemma}[Finite-trace separation]\label{lem:finitetrace}
Let $f_1,\ldots,f_r$ be Laurent polynomials in finitely many variables
over a characteristic-zero field. If
\[
 \sum_{i=1}^r\CT(f_i^N)=0\qquad(N\geq1),
\]
then $\CT(f_i^m)=0$ for every $i$ and $m\geq1$.
\end{lemma}
\begin{proof}
It is enough to consider a number field of coefficients. Indeed, if
one selected moment were nonzero, specialize the finitely generated
$\Q$-algebra of coefficients after inverting that moment. A closed
point gives a number-field specialization preserving it and every
assumed identity.

Fix $m\geq1$ and $1\leq k\leq r$. There are arbitrarily large primes
$p$ splitting completely in the coefficient field and avoiding its
coefficient denominators. For completeness, this infinitude needs no
density theorem. Take an integral primitive element of the Galois
closure. Schur's elementary argument gives infinitely many prime
divisors of values of its minimal polynomial. Outside its discriminant
and index, such a divisor gives a degree-one prime in the Galois field;
all residue degrees are equal, so the prime splits completely.
Schur's argument itself is Euclid's: after translating a nonconstant
integer polynomial to have nonzero constant term $a$, if all prime divisors of its nonzero integer values belonged to a finite
set with product $B$, the polynomial
$f(aBt)/a$ would be integral, nonconstant and congruent to one modulo
$B$. A sufficiently large value has a prime divisor outside that set.

Choose such $p$ larger than the absolute value of every exponent
coordinate in every $f_i^m$, and put
$N=m(1+p+\cdots+p^{k-1})$. In the residue field,
\[
 f_i^N=\prod_{j=0}^{k-1}f_i(z^{p^j})^m.
\]
An exponent equation $\alpha_0+p\alpha_1+\cdots+p^{k-1}\alpha_{k-1}=0$
forces $\alpha_0=0$ by reduction modulo $p$ and the exponent bound.
Repeating gives every $\alpha_j=0$. Hence, writing $a_i=\CT(f_i^m)$,
\[
 \CT(f_i^N)\equiv a_i^k\pmod p.
\]
The assumed trace identity makes $\sum_i a_i^k$ divisible by infinitely
many primes, so it is zero in the number field. We have proved this for
$k=1,\ldots,r$. Newton's identities now give
$\prod_i(T-a_i)=T^r$, so every $a_i=0$. This also contradicts the
nonzero moment preserved in the initial specialization, if one existed.
\end{proof}

Let $K=G^0$ and choose $g_1,\ldots,g_r$ representing $G/(KH)$.
The connected components of $X$ are the affine, open-and-closed orbits
\[
 X_i=K/H_i,\qquad H_i=K\cap g_iHg_i^{-1}.
\]
Consequently $\C[X]=\prod_i\C[X_i]$. Transitivity of $G$ on these
components and normalization imply
\begin{equation}\label{eq:componentaverage}
 \h((f_i))=\frac1r\sum_i\h_i(f_i).
\end{equation}
If $[K,K]\not\subseteq H$, use the connected counterexample on
$K/(K\cap H)$ and extend both functions by zero. This is a regular
operation on an open-and-closed component. Equation~\eqref{eq:componentaverage}
preserves its pure moments and its nonzero marked orders, up to $1/r$.

If $D=[K,K]\subseteq H$, then $D$ lies in every $H_i$, because it is
normal in $G$. Each component is a torus. Embed its character lattice
$M_i=X^*(K/H_i)$ into $M=X^*(K/D)$. For a pure tuple $(f_i)$,
Lemma~\ref{lem:finitetrace} gives purity on every component.
Apply \eqref{eq:dvk} and separate each component support. The maximum
of the finitely many multiplier cutoffs proves safety in
\eqref{eq:componentaverage}. This completes
Theorem~\ref{thm:classification}.

We can also distinguish the Mathieu property from the moment-nullcone
statement. Let $Z_X$ be the pure set and $\mathcal N_G$ the $G$-nullcone:
the functions whose orbit closure contains zero in a finite-dimensional
rational submodule containing the function.

\begin{corollary}\label{cor:nullcone}
One has $Z_X=\mathcal N_G$ if and only if $D\subseteq H$ and
\begin{equation}\label{eq:theta}
 \theta:\bigoplus_i M_i\longrightarrow M,\qquad
 (\chi_i)\longmapsto\sum_i\chi_i
\end{equation}
is injective. The lattices in \eqref{eq:theta} are defined after imposing
$D\subseteq H$.
\end{corollary}
\begin{proof}
The $G$-nullcone equals the $K$-nullcone, since a $G$-orbit closure is
a finite union of translates of the $K$-orbit closure, where $K=G^0$.
For this connected group, the Hilbert--Mumford criterion gives a
one-parameter subgroup with strictly positive weights on any unstable
function \cite[Introduction and Section~2]{Kempf}. These weights
eventually dominate those of any fixed multiplier, so unstable
functions are safe. The non-toral fixed bad
pairs therefore lie outside the nullcone.

In the toral case, finite-trace separation and \eqref{eq:dvk} give a
separate positive functional on every component support. If $\theta$
is injective, these extend to one functional on $M\otimes\Q$, positive
on the entire union of supports. A multiple lifts to a cocharacter of
$K$, since the central torus maps onto $K/D$ with finite kernel. Its
inverse contracts the tuple to zero, with the translation action on
functions. Conversely, a nonzero lattice
relation $\sum_i\chi_i=0$ gives a pure tuple with component $\chi_i$
when $\chi_i\ne0$ and zero otherwise. To exclude zero from its orbit
closure, let $I=\{i:\chi_i\ne0\}$ and let $v_i\in\C[X]$ equal
$\chi_i$ on $X_i$ and zero elsewhere. The tuple is $v=\sum_{i\in I}v_i$
in the finite-dimensional rational $K$-module
$W=\bigoplus_{i\in I}\C v_i$. If $z_i$ are the coordinate functions
on $W$, then $\prod_{i\in I}z_i\in\C[W]^K$ because
$\sum_i\chi_i=0$. This polynomial has value one at $v$ and zero at
the origin. Thus $v$ is outside the nullcone. The product here is a
polynomial on $W$, rather than a product of component-supported
functions in $\C[X]$.
Finally, all nullcone elements are pure because the moments are invariant
homogeneous polynomials vanishing at zero.
\end{proof}

\section{Recompositions and ordinary Cauchy transforms}\label{sec:recompositions}

Return to either exceptional pair $P,Q$ in
Theorem~\ref{thm:exceptional}, with $\Gamma=2O$ or $2I$, respectively.
Write
\[
 (c,a)=(28,3)\text{ or }(66,5),\qquad b_j=\h(QP^j).
\]
The facts needed in this section are
\begin{equation}\label{eq:denominatorfacts}
 b_0=0,\qquad
 \operatorname{den}(b_j)\mid(cj-a)!\ (j\geq1),\qquad
 v_p(b_j)=-1\text{ when }p=cj-a\text{ is prime}.
\end{equation}
The first is a weight calculation, the second is
Proposition~\ref{prop:markedcomponent}, and the third is
Theorem~\ref{thm:exceptional}. Here $\operatorname{den}(b_j)$ is the
positive denominator of the reduced rational number $b_j$.

For regular $S,V$ on $\SL_2/\Gamma$, evaluate the Haar integral on the
compact quotient $\SU(2)/\Gamma$. The ordinary marked Cauchy transform
is the convergent germ
\[
 C_V^S(z)=\h\left(\frac{V}{z-S}\right)
 =\sum_{m\geq0}\frac{\h(VS^m)}{z^{m+1}}
 \quad\text{for }|z|>\max|S|.
\]
We use this ordinary series throughout this section.

\begin{theorem}\label{thm:recompositions}
For each $d\geq1$, choose an integer $e$ satisfying
\begin{equation}\label{eq:choosemarker}
 0\leq e<d,\qquad \gcd(cd,ce-a)=1,
\end{equation}
and set $Q_d=QP^e$. Such an $e$ always exists. For every complex
polynomial $H$ of degree $d$, the same multiplier $Q_d$ satisfies
\[
 \h\bigl(Q_dH(P)^m\bigr)\ne0\quad\text{for infinitely many }m,
 \qquad C_{Q_d}^{H(P)}(z)\text{ is transcendental over }\C(z).
\]
Furthermore
\begin{equation}\label{eq:unitcomposed}
 \h\bigl(H(P)^m\bigr)=H(0)^m,\qquad
 C_1^{H(P)}(z)=\frac1{z-H(0)}.
\end{equation}
In particular, every nonzero polynomial recomposition with $H(0)=0$
is pure and unsafe.
\end{theorem}

The multiplier depends on the degree, not on the coefficients of $H$.
For example, the smallest possible $e$ in
\eqref{eq:choosemarker} gives an unambiguous choice.

We first record two elementary algebraic facts used to allow arbitrary
complex coefficients.

\begin{lemma}\label{lem:goodvaluations}
Let $L$ be a finitely generated field over $\Q$, and let $S\subset L^*$
be finite. For all but finitely many rational primes $p$, there is a
valuation $w$ on $L$ extending $v_p$, normalized by $w(p)=1$, for which
every member of $S$ is a unit.
\end{lemma}
\begin{proof}
Choose a transcendence basis $t$ with $L$ finite over $\Q(t)$.
On $\Q(t)$ use the $p$-adic Gauss valuation: on a polynomial it is
the minimum coefficient valuation, and on a quotient it is the difference.
Choose monic equations over $\Q(t)$ for each $\alpha\in S$ and its
inverse. All coefficients of these finitely many equations are integral
for all but finitely many Gauss valuations. Extend the valuation to $L$;
existence follows from the valuation-ring domination lemma
\cite[Lemma~10.50.2]{Stacks}. The monic equations force both
$w(\alpha)\geq0$ and $w(\alpha^{-1})\geq0$, so $w(\alpha)=0$.
Values in an algebraic extension are rational multiples of $w(p)$:
in an algebraic equation the minimum valuation must occur at least
twice, which gives an integer multiple of the element's value in
the original value group. This justifies the normalization.
\end{proof}

\Needspace{9\baselineskip}
\begin{lemma}[Finite-generation form of Eisenstein's argument]
\label{lem:eisenstein}
If $L$ has characteristic zero and $A(t)\in L[[t]]$ is algebraic over
$L(t)$, all coefficients of $A$ belong to one finitely generated
$\Z$-subalgebra of $L$.
\end{lemma}
\begin{proof}
Choose $F(t,Y)\in L[t,Y]$ with $F(t,A)=0$ and $F_Y(t,A)\ne0$;
characteristic zero guarantees such a relation. Put
$r=\ord_tF_Y(t,A)$, choose $N>r$, and write
$A=a+t^NB$, where $a=A\bmod t^N$. Taylor expansion gives
\[
 G(t,Y)=t^{-N-r}F(t,a+t^NY)\in L[t,Y].
\]
For the constant term, divisibility follows by comparing with $F(t,A)=0$;
all terms of order at least two have $t$-order at least $2N>N+r$.
The coefficient of $Y$ has a nonzero constant term $u$, and all higher
powers of $Y$ have zero constant term. Thus $G(0,Y)$ is linear.
Successive coefficient comparison in $G(t,B)=0$ uses only addition,
multiplication and division by $u$. It puts all coefficients of $B$
in the ring generated by the coefficients of $G$ and $u^{-1}$.
Adjoining the finitely many coefficients of $a$ proves the lemma.
This is the part of Eisenstein's theorem needed here; see
\cite{Rond} for its classical and more general forms.
\end{proof}

\begin{proof}[Proof of Theorem~\ref{thm:recompositions}]
Since $\gcd(c,a)=1$, the choice in \eqref{eq:choosemarker} is possible.
For every prime $\ell\mid d$ with $\ell\nmid c$, choose
$e\not\equiv a/c\pmod\ell$. A prime dividing $c$ cannot divide $a$,
so imposes no condition. The Chinese remainder theorem combines the
choices modulo $\operatorname{rad}(d)$, and a representative lies below $d$.

Write $H(T)=\sum_{j=0}^d h_jT^j$, and let $L$ be the field generated by
its coefficients over $\Q$. For a sufficiently large prime
\begin{equation}\label{eq:composedprime}
 p=c(dm+e)-a,
\end{equation}
Lemma~\ref{lem:goodvaluations} gives a $w$ making every nonzero $h_j$
a unit. Expansion of $H(T)^m$ yields
\[
 \h\bigl(Q_dH(P)^m\bigr)
 =h_d^m b_{dm+e}
   +\sum_{j<dm}[T^j]H(T)^m\,b_{j+e}.
\]
The first summand has valuation $-1$ by
\eqref{eq:denominatorfacts}. Every other summand is integral: its
coefficients are integral, and $c(j+e)-a<p$ makes its factorial
denominator prime to $p$. The term $b_0$ is zero. A unique term of
negative valuation cannot cancel. The progression
\eqref{eq:composedprime} is coprime by \eqref{eq:choosemarker}; Dirichlet's
theorem supplies infinitely many such primes after any finite exclusions.
This proves the fixed-multiplier nonvanishing assertion.

Suppose the ordinary series were algebraic over $\C(z)$.
It would be algebraic over $L(z)$ as well. To see this descent, a
polynomial relation with fixed degree bounds is a linear system over
$L$ in finitely many unknown coefficients. Extending scalars does not
change its rank. Even though the system has infinitely many rows, its
row space has a finite basis, so a nonzero complex solution gives a
nonzero $L$ solution.

Apply Lemma~\ref{lem:eisenstein} to
$A(t)=\sum_{m\geq0}\h(Q_dH(P)^m)t^m$. There is a finite set
$S_A\subset L^*$ whose $\Z$-algebra contains every coefficient of $A$.
For a sufficiently large prime in \eqref{eq:composedprime}, apply
Lemma~\ref{lem:goodvaluations} simultaneously to $S_A$ and the nonzero
coefficients of $H$. It makes every coefficient of $A$ integral, whereas
the preceding expansion gives valuation $-1$ for its $m$th coefficient.
This contradiction proves transcendence.
Finally, all positive powers of $P$ integrate to zero, so expanding
$H(P)^m$ leaves only $H(0)^m$. Summing its geometric series proves
\eqref{eq:unitcomposed}.
\end{proof}

Thus the unit Cauchy transform is rational while the marked transform
is transcendental for the specified fixed polynomial multiplier.

\subsection*{Coefficient algebras}

There is also an exact statement with nilpotents retained. Let $R$ be
any commutative unital $\C$-algebra, and extend $\h$ $R$-linearly to
$\h_R:R\otimes_{\C}\C[\SL_2/\Gamma]\to R$. An element is
\emph{eventually pure} if its positive moments vanish for all sufficiently
large orders. Safety in the following statement is against every
multiplier in this entire algebra.

\begin{corollary}\label{cor:coefficientrings}
For $H(T)=\sum_{j=0}^d c_jT^j\in R[T]$,
\[
 \begin{aligned}
 H(P)\text{ is pure}&\ \Longleftrightarrow\ c_0=0,\\
 H(P)\text{ is eventually pure}&\ \Longleftrightarrow\ c_0
          \text{ is nilpotent},\\
 H(P)\text{ is safe}&\ \Longleftrightarrow\ \text{every }c_j
          \text{ is nilpotent}
 \ \Longleftrightarrow\ H(P)\text{ is nilpotent}.
 \end{aligned}
\]
\end{corollary}
\begin{proof}
The formula $\h_R(H(P)^m)=c_0^m$ proves the first two assertions.
If every coefficient is nilpotent, their finitely generated ideal is
nilpotent. Explicitly, $c_j^{r_j}=0$ makes its
$1+\sum_j(r_j-1)$ power zero. Thus $H(P)$ is nilpotent and safe.

Conversely, a nonnilpotent coefficient survives in a prime quotient
of $R$. Pass to the fraction field of that quotient. If the image of
$H$ has positive degree, the valuation argument in
Theorem~\ref{thm:recompositions} applies to the finitely generated
subfield over $\Q$ generated by its coefficients. It gives a fixed
multiplier $QP^e$ with infinitely many nonzero
moments. Those moments were therefore nonzero in $R$. If the image of
$H$ is a nonzero constant, the multiplier one already detects failure.
Hence safety forces every coefficient to be nilpotent. Since any
nilpotent function is safe, the final equivalence follows too.
\end{proof}

\appendix
\section{Invariance of the exceptional binary forms}\label{app:invariance}

We give a verification using the root sets; it also fixes the embeddings
of the binary groups. Identify $\mathbb P^1$ with the unit sphere by
\[
 q\longmapsto\frac{(2\Re q,2\Im q,|q|^2-1)}{1+|q|^2}.
\]
The projective action of $\SU(2)$ is the rotation action.

For $f_8$, the finite roots satisfy $q^4=-7\pm4\sqrt3$.
They have arguments $(2k+1)\pi/4$ and squared moduli $2\pm\sqrt3$.
Their sphere coordinates are exactly
$(\pm1,\pm1,\pm1)/\sqrt3$, the eight vertices of a cube.
Its rotation group is $S_4$, and its inverse image in $\SU(2)$ is $2O$.
The root set makes the line $\C f_8$ stable. Its scalar character factors
through $S_4$, because the central element $-I$ acts trivially in even
degree. The matrix
$\operatorname{diag}(e^{\pi i/4},e^{-\pi i/4})$ fixes $f_8$ directly
and projects to a quarter-turn, a four-cycle on the body diagonals.
Every character of $S_4$ has the same value $\epsilon\in\{1,-1\}$
on its transpositions and value $\epsilon^3=\epsilon$ on a four-cycle.
Thus this character is trivial and $f_8$ itself is invariant.

For $f_{12}$, let $\varphi=(1+\sqrt5)/2$. Besides $0$ and $\infty$,
its roots are
\[
 q=\varphi^{-1}e^{2\pi i k/5},\qquad
 q=\varphi e^{(2k+1)\pi i/5}\quad(0\leq k<5).
\]
Their sphere images are two poles and two regular pentagons at heights
$-1/\sqrt5$ and $1/\sqrt5$, of radius $2/\sqrt5$, offset by $\pi/5$.
These are the vertices of a regular icosahedron. Its rotation group is
$A_5$ and its inverse image is $2I$. Again the line spanned by the
binary form is stable, and its scalar character factors through $A_5$.
A three-cycle is conjugate to its inverse in $A_5$, so its character
value has both square and cube equal to one. It is therefore one.
Three-cycles generate $A_5$, proving invariance.

\section{The differential extraction rule}\label{app:cartier}

Here we prove the coordinate independence used in
Section~\ref{sec:klein}. Let $L$ be a one-variable function field over
$\Fp$, and choose a separating element $t$. Then
$L=L^p(t)$ and $1,t,\ldots,t^{p-1}$ is a basis over $L^p$.
Indeed, $L/L^p(t)$ is purely inseparable but is also separable, as a
subextension of the separable extension $L/\Fp(t)$. It is therefore
trivial; and $t\notin L^p$, since $\dd t\ne0$.

Every differential is consequently $\sum_{i=0}^{p-1}a_i^pt^i\dd t$.
Modulo exact differentials, only the last term remains: differentiating
$\sum b_i^pt^i$ shows that the exact differentials have the other
$p-1$ basis terms. Thus $\Omega_L/\dd L$ is one-dimensional over $L^p$,
with basis $[t^{p-1}\dd t]$.

For $u\in L$ put $\gamma(u)=[u^{p-1}\dd u]$ in this quotient.
It satisfies
\[
 \gamma(uv)=u^p\gamma(v)+v^p\gamma(u),\qquad
 \gamma(u+v)=\gamma(u)+\gamma(v).
\]
The product formula is literal. For additivity, differentiate the integer
polynomial $((U+V)^p-U^p-V^p)/p$ and reduce modulo $p$; its differential
is exactly the difference of the two sides before taking classes.
These rules also imply $\gamma(u^p)=0$.

Writing $u=\sum a_i^pt^i$ now gives
\[
 [u^{p-1}\dd u]=\left(\frac{\dd u}{\dd t}\right)^p
                         [t^{p-1}\dd t].
\]
If $u$ is another separating coordinate, the class of a differential
written as $a^{p}u^{p-1}\dd u$ is therefore
$(a\,\dd u/\dd t)^p[t^{p-1}\dd t]$.
The extracted differential is $a\dd u$ in either coordinate.
This proves that the rule defining $\Cart$ is well defined.
For a separable extension of function fields, the same $t$ is separating
in both and the same expansion works in both, proving compatibility
with pullback. Formula~\eqref{eq:cartierrule} follows by grouping powers
whose exponents are congruent to $p-1$. See also \cite[Section~2.2]{AchterHowe} for
the usual Cartier--Manin terminology.

\section{Checking the polynomial identities}\label{app:identities}

Here are explicit forms of the algebra in
\eqref{eq:kleinrelations}, so that the proof does not require a table
of invariant-theoretic normalizations. For $n=8$, put
\[
 B=XY(X^4-Y^4),\qquad
 T=X^{12}-33X^8Y^4-33X^4Y^8+Y^{12}.
\]
Differentiation gives $G=B^2$ and $J=BT/6$. To check
$T^2=f_8^3-108B^4$, set $s=X^4+Y^4$ and $u=X^4Y^4$.
Then $f_8=s^2+12u$, $T=s(s^2-36u)$ and $B^4=u(s^2-4u)^2$.
The identity follows by expanding these three short expressions.

For $n=12$, differentiation gives $12J=-T$, where
\[
 T=X^{30}+Y^{30}+522(X^{25}Y^5-X^5Y^{25})
                  -10005(X^{20}Y^{10}+X^{10}Y^{20}).
\]
With $t=X^5/Y^5$, the required identity becomes the following
one-variable polynomial identity:
\begin{align*}
 &(t^6+522t^5-10005t^4-10005t^2-522t+1)^2\\
 &\qquad=(t^4-228t^3+494t^2+228t+1)^3
              +1728t(t^2+11t-1)^5.
\end{align*}
Multiplying the displayed polynomials verifies it. Homogenization then
gives $T^2=G^3+1728f_{12}^5$, hence
$J^2=G^3/144+12f_{12}^5$ as asserted.

\paragraph{AI assistance and review.}
ChatGPT (OpenAI) assisted with mathematical arguments, literature checks,
exposition, proof review and verification code.

\begin{thebibliography}{10}

\bibitem{AchterHowe}
J.~D.~Achter and E.~W.~Howe,
\emph{Hasse--Witt and Cartier--Manin matrices: a warning and a request},
\href{https://arxiv.org/abs/1710.10726v5}{arXiv:1710.10726v5},
corrected preprint, 2020.

\bibitem{AdamsTaibi}
J.~Adams and O.~Ta\"ibi, \emph{Galois and Cartan cohomology of real groups},
preprint, 2018, Section~3.
\href{https://otaibi.perso.math.cnrs.fr/galois.pdf}{Author's manuscript}.

\bibitem{Arzhantsev}
I.~V.~Arzhantsev, \emph{Invariant ideals and Matsushima's criterion},
\href{https://arxiv.org/abs/math/0506430v1}{arXiv:math/0506430v1}, 2005.

\bibitem{Borel}
A.~Borel, \emph{Linear Algebraic Groups}, second edition,
Graduate Texts in Mathematics 126, Springer, 1991.

\bibitem{DvK}
J.~J.~Duistermaat and W.~van der Kallen,
\emph{Constant terms in powers of a Laurent polynomial},
Indag. Math. (N.S.) \textbf{9} (1998), 221--231.
\href{https://doi.org/10.1016/S0019-3577(98)80020-7}{doi:10.1016/S0019-3577(98)80020-7}.

\bibitem{Flicker}
Y.~Z.~Flicker, \emph{Finite subgroups of SL(2,F) and automorphy},
Manuscripta Math. \textbf{167} (2022), 213--243.
\href{https://doi.org/10.1007/s00229-021-01274-x}{doi:10.1007/s00229-021-01274-x}.

\bibitem{Kempf}
G.~Kempf, \emph{Instability in invariant theory}, original preprint,
transcription by I.~Morrison,
\href{https://arxiv.org/abs/1807.02890v1}{arXiv:1807.02890v1}, 2018.

\bibitem{LongSU2}
C.~D.~Long, \emph{Counterexamples to the $xz$-Conjecture and the
Mathieu Conjecture for $\SU(2)$},
\href{https://arxiv.org/abs/2607.19012v1}{arXiv:2607.19012v1}, 2026.

\bibitem{LongCompact}
C.~D.~Long, \emph{The Mathieu Property for Compact Connected Lie Groups},
preprint, 2026,
\href{https://arxiv.org/abs/2609.16178v1}{arXiv:2609.16178v1}.

\bibitem{LongSpherical}
C.~D.~Long, \emph{A Classification of Mathieu's Second Conjecture for
Reductive Spherical Pairs}, September 2026.
\href{https://github.com/long-mathematics/mathieu-second-conjecture-spherical-pairs/blob/e9edac262b4b64eb68feb0e6e80496cefe4a1dee/mathieu_second_conjecture_spherical_pairs.tex}{Versioned GitHub source, revision e9edac2}.

\bibitem{Mathieu}
O.~Mathieu, \emph{Some conjectures about invariant theory and their
applications}, S\'eminaires et Congr\`es \textbf{2} (1997), 263--279.
\href{https://smf.emath.fr/sites/default/files/2026-04/smf_sem-cong_2_263-279__sample.pdf}{Publisher's excerpt, pp.~263--267}.

\bibitem{Molev}
A.~I.~Molev, \emph{Gelfand--Tsetlin bases for classical Lie algebras},
in Handbook of Algebra, vol.~4, Elsevier, 2006, 109--170.
\href{https://arxiv.org/abs/math/0211289v2}{arXiv:math/0211289v2}.

\bibitem{Rond}
G.~Rond, \emph{About Eisenstein's theorem},
\href{https://arxiv.org/abs/2412.11738v2}{arXiv:2412.11738v2}, 2024.

\bibitem{Stacks}
The Stacks Project Authors, \emph{The Stacks Project},
\href{https://stacks.math.columbia.edu/tag/00I8}{Tag 00I8},
Lemma~10.50.2.

\bibitem{Sutherland}
A.~V.~Sutherland, \emph{Dirichlet $L$-functions, primes in arithmetic
progressions}, MIT 18.785, Lecture~18, 2021, Theorem~18.1.
\href{https://ocw.mit.edu/courses/18-785-number-theory-i-fall-2021/mit18_785f21_lec18.pdf}{Lecture notes}.

\end{thebibliography}
\end{document}
